\documentclass[12pt,a4paper]{article}
\usepackage[utf8]{inputenc}
\usepackage[T1]{fontenc}
\usepackage[french]{babel}
\usepackage{geometry}
\usepackage{longtable}
\usepackage{booktabs}
\usepackage{array}
\usepackage{hyperref}
\usepackage{xcolor}
\usepackage{enumitem}
\usepackage{titlesec}

\geometry{margin=2.2cm}
\setlist[itemize]{leftmargin=1.2cm}
\setlist[enumerate]{leftmargin=1.1cm}
\hypersetup{colorlinks=true,linkcolor=black,urlcolor=blue}
\titleformat{\section}{\large\bfseries}{\thesection.}{0.6em}{}
\titleformat{\subsection}{\normalsize\bfseries}{\thesubsection.}{0.6em}{}

\title{\textbf{Rapport d'Audit Technique}\\Smart Shepherd - Plateforme Admin IoT/IA}
\author{Audit architecture, produit et ingenierie}
\date{29 mai 2026}

\begin{document}
\maketitle
\tableofcontents
\newpage

\section{Resume executif}
Ce rapport presente une analyse technique complete du dashboard administratif et de la plateforme Smart Shepherd sur la base du code source disponible.\

\textbf{Conclusion courte :} la plateforme est avancee et convaincante en demonstration (temps reel, carte, analytics, IA, simulation), mais reste hybride entre MVP avance et produit SaaS industriel.\

\textbf{Forces majeures :}
\begin{itemize}
  \item Stack moderne coherente pour un produit IoT (React + Node + MongoDB + MQTT + Socket.IO + FastAPI).
  \item Dashboard riche avec nombreuses pages metier et UX globalement solide.
  \item Simulation IoT detaillee et parametres avances de scenario.
  \item Presence d'une base DevOps (Docker Compose, Prometheus, Grafana, Alertmanager).
\end{itemize}

\textbf{Risques majeurs :}
\begin{itemize}
  \item Securite et authentification heterogenes (plusieurs middlewares auth, bypass dev, roles non uniformes).
  \item IA multi-chemin (Node local, FastAPI, LLM externe) avec fallback mock partiellement opaque.
  \item Contrats API et evenements temps reel non totalement uniformises.
  \item Presence de comportements fallback/mocks pouvant masquer des incidents en production.
\end{itemize}

\section{Methodologie et perimetre}
\subsection{Sources inspectees}
L'analyse couvre principalement :
\begin{itemize}
  \item Frontend: routage, pages, hooks et services temps reel.
  \item Backend: API Express, routes metier, MQTT, Socket.IO, auth.
  \item Microservice IA Python FastAPI.
  \item Infrastructure Docker (dev/prod) et observabilite.
\end{itemize}

\subsection{Limites}
\begin{itemize}
  \item Pas de capture d'ecran exploitee dans ce passage.
  \item Pas de benchmark de charge execute dans ce rapport.
  \item Evaluation basee sur le code et la structure projet.
\end{itemize}

\section{1. Identification des technologies}
\subsection{Frontend}
\begin{itemize}
  \item React 18 + TypeScript + Vite.
  \item React Router v6 pour les routes applicatives.
  \item Zustand pour l'etat global IoT et alertes.
  \item React Query pour requetes API et mutations.
  \item Tailwind CSS + composants UI custom.
  \item Chart.js/Recharts pour visualisation.
  \item Leaflet/React-Leaflet + Geoman pour cartes/geofence.
  \item MQTT client navigateur et Socket.IO client.
\end{itemize}

\subsection{Backend}
\begin{itemize}
  \item Node.js + Express (ESM).
  \item MongoDB via Mongoose.
  \item Redis (cache conditionnel).
  \item Socket.IO serveur pour push realtime.
  \item MQTT client serveur pour ingestion/emission flux IoT.
\end{itemize}

\subsection{IA et data science}
\begin{itemize}
  \item Service IA Node (tfjs-node conditionnel, logique prediction et scoring).
  \item Microservice FastAPI (TensorFlow/XGBoost, endpoints IA).
  \item LLM externe Anthropic avec fallback local heuristique.
\end{itemize}

\subsection{Infra et exploitation}
\begin{itemize}
  \item Docker Compose dev/prod.
  \item Mosquitto, Mongo, Redis.
  \item Prometheus, Alertmanager, Grafana.
  \item Nginx reverse proxy.
\end{itemize}

\subsection{LoRa / LoRaWAN}
\begin{itemize}
  \item Presence d'un dossier firmware ESP32 avec bibliotheques LoRaWAN.
  \item Pas d'integration backend explicite complete type network server LoRaWAN -> decoder -> pipeline data.
\end{itemize}

\textbf{Verdict coherence techno :} choix global logique et pertinent pour une plateforme IoT intelligente; les points faibles proviennent surtout de l'uniformisation d'architecture et de gouvernance runtime.

\section{2. Analyse complete des pages}
\subsection{Pages identifiees}
\begin{itemize}
  \item Dashboard, Map, Animals, Animal Profile, Compare, Agenda, Alerts.
  \item Anomalies, Analytics, AI Dashboard.
  \item Users, Hardware, Settings.
  \item Admin: AI Settings, Labelling.
\end{itemize}

\subsection{Evaluation fonctionnelle par domaine}
\textbf{Operations:}
\begin{itemize}
  \item Dashboard: KPIs, widgets realtime, alertes, mini-map, meteo.
  \item Map: supervision geospatiale, geofencing, edition de zones.
  \item Animals: liste virtualisee, tri/filtre, score sante, navigation detail.
  \item Alerts: centre de triage realtime.
  \item Agenda: planification vet.
\end{itemize}

\textbf{Intelligence:}
\begin{itemize}
  \item Analytics: heatmap, batterie, tendances historiques, export.
  \item Anomalies: registre + analyse IA asynchrone.
  \item AI Dashboard: vue predictive, severite, recommandations.
\end{itemize}

\textbf{Administration:}
\begin{itemize}
  \item Users: CRUD avec optimistic update (et mode mock selon config).
  \item Hardware: flotte colliers, etat derive du store.
  \item Settings/AI Settings: configuration simulation, seuils, canaux notifications.
  \item Labelling: annotation terrain alimentee par les alertes et le troupeau live, avec fallback store en secours.
\end{itemize}

\subsection{UX/UI}
\begin{itemize}
  \item Niveau visuel global bon a tres bon (hierarchie, composants, skeletons, badges severite).
  \item Navigation claire via piliers metier.
  \item Risque de surcharge informationnelle sur certaines vues denses.
  \item Certains fallbacks UX peuvent masquer des incidents techniques reels.
\end{itemize}

\section{3. Analyse speciale IA}
\subsection{Constat de realite IA}
La plateforme propose une vraie couche IA, mais elle fonctionne selon plusieurs modes selon la disponibilite des dependances, des modeles et des cles API.

\subsection{Modes detectes}
\begin{itemize}
  \item IA locale Node (prediction/scoring).
  \item IA microservice FastAPI (health-check, battery-rul, position).
  \item IA generative externe (Anthropic) pour synthese.
  \item Fallback heuristique locale en cas d'echec API/cle absente, expose via un mode IA explicite.
\end{itemize}

\subsection{Points critiques IA}
\begin{itemize}
  \item Risque de confusion produit fortement reduit: le backend expose maintenant un \texttt{model\_mode} explicite (anthropic/local/local-cache).
  \item Endpoints avec retour success/fallback toujours a surveiller, mais la transparence d'incident est meilleure qu'avant.
  \item Absence d'une gouvernance MLOps complete (versioning modele, quality gates, drift monitoring).
\end{itemize}

\subsection{Recommandations IA prioritaires}
\begin{enumerate}
  \item Standardiser la reponse IA avec un champ obligatoire \texttt{model\_mode} (real/mock/fallback).
  \item Centraliser l'inference via FastAPI; garder Node en orchestrateur et policy layer.
  \item Mettre en place registre de modeles + metriques qualite + seuils d'acceptation.
  \item Isoler IA explicative (LLM) de l'IA decisionnelle (scoring risque).
\end{enumerate}

\section{4. Analyse de la simulation}
\subsection{Etat actuel}
\begin{itemize}
  \item Simulation riche (nombre animaux, vitesse, densite, scenarios, batterie, alertes).
  \item Simulation encore principalement centree frontend/store, avec un pipeline temps reel partiellement unifie.
  \item Bonne valeur pour demo, UX tests, essais fonctionnels rapides.
\end{itemize}

\subsection{Limites}
\begin{itemize}
  \item Pas de digital twin backend canonique.
  \item Realisme biologique/comportemental limite par des regles heuristiques.
  \item Rejouabilite et tracabilite simulation encore perfectibles.
\end{itemize}

\subsection{Cible recommandee}
\begin{enumerate}
  \item Moteur de simulation backend event-driven.
  \item Publication simulation via MQTT sur le meme pipeline que le reel.
  \item Scenarios parametrables avec \texttt{run\_id} et seeds reproductibles.
  \item Correlation simulation avec meteo, geofence, pannes capteurs et anomalies IA.
\end{enumerate}

\section{5. Analyse du dashboard administratif}
\subsection{Architecture frontend}
\begin{itemize}
  \item Organisation modulaire globalement bonne (pages, composants, hooks).
  \item Etat global coherent via Zustand, mais certaines logiques restent dupliquees (MQTT/realtime/services).
  \item Presence d'une forte couche de presentation avancee.
\end{itemize}

\subsection{Roles et permissions}
\begin{itemize}
  \item Intention RBAC visible cote UI.
  \item Incoherences role naming detectees (admin/ADMIN/super\_admin selon couches).
  \item La securite ne doit pas dependre de la seule couche UI.
\end{itemize}

\subsection{Securite}
\begin{itemize}
  \item Corrections recentes sur JWT\_SECRET positives.
  \item Mais coexistence de plusieurs middlewares auth et bypass dev a controler strictement.
\end{itemize}

\section{6. Detection des problemes}
\subsection{Problemes critiques}
\begin{enumerate}
  \item Heterogeneite des middlewares d'authentification et des conventions roles.
  \item Bypass dev potentiellement actif selon environnement et routes.
  \item Evenements realtime non completement alignes entre emit backend et ecoute frontend.
\end{enumerate}

\subsection{Problemes majeurs}
\begin{enumerate}
  \item Contrats API non uniformes (formes de payload variees).
  \item Multiples fallback/mocks en chemins utilisateur sans indicateur unifie.
  \item Dette de standardisation du pipeline IA.
\end{enumerate}

\subsection{Problemes moyens}
\begin{enumerate}
  \item Couplage fort de certains composants pages/services.
  \item Presence de logiques historiques/optionnelles qui complexifient la maintenance.
  \item Documentation fonctionnelle riche mais parfois en decalage avec execution runtime.
\end{enumerate}

\section{7. Fonctionnalites manquantes pour un niveau SaaS IoT pro}
\begin{itemize}
  \item RBAC centralise avec policy engine unique.
  \item Audit logs complets (auth, actions admin, config, IA).
  \item Device management industriel: provisioning, rotation credentials, OTA, inventaire firmware reel.
  \item Observabilite avancee: traces distribuees, SLO, runbooks.
  \item MLOps: monitoring de drift, performance modele, dataset lineage.
  \item Gouvernance multi-tenant stricte et retention politique.
  \item Exports reglementaires et mode degrade explicite.
\end{itemize}

\section{8. Analyse architecture logicielle}
\subsection{Evaluation}
\begin{itemize}
  \item Modularite frontend: bonne.
  \item Backend: monolithe enrichi avec extensions; viable mais a cadrer.
  \item Event-driven: present mais non strictement unifie.
  \item IA microservice: bonne direction, integration encore hybride.
\end{itemize}

\subsection{Architecture cible recommandee}
\begin{enumerate}
  \item API Gateway + Auth service unique + RBAC policy centralisee.
  \item Ingestion IoT unique (MQTT) + normalisation schema + event bus interne.
  \item Service IA central d'inference versionnee.
  \item Contrats API/evenements versionnes (OpenAPI + event schema).
  \item Frontend branche sur etats data explicites (live/simulated/cached/mock).
\end{enumerate}

\subsection{Architecture cible enterprise}
\begin{itemize}
  \item Couche d'entree: API Gateway, Auth Service, rate limiting et policies de securite.
  \item Couche metier: IoT Ingestion Service, Alert Service, Analytics Service, Device Service.
  \item Couche IA: FastAPI centralise, model registry, inference service, monitoring des modeles.
  \item Couche temps reel: MQTT Broker, normalizer, event bus, WebSocket Gateway, subscriptions front.
  \item Couche donnees: MongoDB pour les objets metier, PostgreSQL pour les registres structurants, Redis pour cache/queues.
  \item Couche front: dashboard modulaire, design system, state management explicite live/sim/mock/fallback.
\end{itemize}

\subsection{Contrats, schemas et standardisation}
\begin{itemize}
  \item Toutes les APIs doivent exposer \texttt{\{success, data, error, metadata\}}.
  \item Les erreurs doivent etre normalisees par code, message, details et correlation\_id.
  \item Les evenements MQTT et WebSocket doivent partager le meme schema canonique.
  \item Les DTO doivent etre versionnes (v1, v2) avec validation Zod/Joi cote backend.
  \item OpenAPI doit servir de source de verite pour les routes HTTP.
\end{itemize}

\subsection{Architecture temps reel cible}
\begin{enumerate}
  \item IoT Device publie vers un topic MQTT normalise.
  \item Broker MQTT alimente un event normalizer avec idempotence et correlation id.
  \item Event bus distribue les evenements vers alerting, analytics, IA et websocket.
  \item WebSocket Gateway pousse vers le dashboard uniquement les vues abonnees.
  \item Le frontend affiche clairement les modes LIVE, SIMULATION, MOCK et FALLBACK.
\end{enumerate}

\subsection{IA, MLOps et gouvernance}
\begin{itemize}
  \item Un service IA decisionnel doit produire health score, anomaly score et behavior analysis.
  \item Un service IA explicatif doit generer les syntheses textuelles et recommandations.
  \item Le model registry doit conserver version, dataset, metrics, date de promotion et rollback target.
  \item Les predictions doivent remonter un \texttt{confidence score} et un \texttt{model\_mode}.
  \item Les modeles doivent etre surveilles sur drift, calibration, latency et taux d'erreur.
\end{itemize}

\subsection{Simulation enterprise}
\begin{itemize}
  \item La simulation doit devenir un moteur backend event-driven avec \texttt{run\_id} et seeds reproductibles.
  \item Chaque scenario doit simuler GPS, temperature, rythme cardiaque, batterie, geofence et panne reseau.
  \item Les evenements de simulation doivent passer par le meme pipeline MQTT que les donnees reales.
  \item Le dashboard doit montrer explicitement la provenance des donnees: LIVE, SIMULATION, MOCK ou FALLBACK.
\end{itemize}

\subsection{Observabilite et DevOps}
\begin{itemize}
  \item Metriques: latence API, latence MQTT, files d'attente, erreurs IA, taux d'alertes, uptime.
  \item Logs: structure JSON, correlation\_id, audit logs admin, logs d'inference IA.
  \item Traces: OpenTelemetry pour router les requetes entre gateway, services et broker.
  \item Alerting: Prometheus + Grafana + Alertmanager + notifications critiques.
  \item Protection: backups, rollback, health checks, readiness/liveness et rotation secrets.
\end{itemize}

\subsection{Device management}
\begin{itemize}
  \item Inventaire firmware et hardware par device.
  \item Provisioning, rotation de certificats, stockage des identites de devices.
  \item Heartbeat, batterie, connectivite, geofence, statut de securite.
  \item OTA updates avec suivi d'etat et rollback firmware.
\end{itemize}

\subsection{Migration strategy}
\begin{enumerate}
  \item Phase A: introduire Auth Service, schema API commun et modele de reponse standard.
  \item Phase B: unifier les evenements realtime et isoler le pipeline IA.
  \item Phase C: deplacer la simulation vers le backend et brancher le dashboard sur les sources canoniquees.
  \item Phase D: industrialiser observabilite, MLOps et device management.
\end{enumerate}

\subsection{Backlog critique priorise}
\begin{longtable}{p{4.2cm}p{1.8cm}p{8.4cm}}
\toprule
\textbf{Chantier} & \textbf{Priorite} & \textbf{Effet attendu} \\
\midrule
Auth/RBAC centralise & P0 & Une seule source d'autorite pour les roles, permissions et tokens \\
API/event contracts & P0 & Reduction des incoherences de payload et des integrations fragiles \\
Realtime unifie & P0 & Moins de duplications et de disconnects inter-pages \\
IA centralisee & P1 & Modele explicite, traçable et monitorable \\
Simulation backend & P1 & Donnees reproductibles et pipeline canonique \\
Observabilite & P1 & Diagnostic rapide des incidents et mesure SLA/SLO \\
Device management & P2 & Ready pour exploitation terrain serieuse \\
CI/CD et security hardening & P2 & Livraison plus sure et plus frequentable \\
\bottomrule
\end{longtable}

\section{9. Evaluation professionnelle (notes /10)}
\begin{longtable}{p{8.5cm}p{2cm}p{4cm}}
\toprule
\textbf{Axe} & \textbf{Note} & \textbf{Commentaire court} \\
\midrule
Architecture logicielle & 7.2 & Bonne base, besoin d'unification \\
Frontend admin & 8.4 & Niveau visuel et UX eleve \\
Backend API & 6.8 & Fonctionnel mais heterogene \\
IA (maturite produit) & 6.1 & Potentiel fort, gouvernance insuffisante \\
UX/UI & 8.2 & Solide et moderne \\
Temps reel IoT & 7.0 & Bon socle, schemas a aligner \\
Securite applicative & 5.9 & Chantiers prioritaires restants \\
Qualite produit SaaS & 7.0 & MVP avance, pas encore enterprise \\
Credibilite startup tech demo & 8.0 & Convaincant pour demo/invest \\
Niveau ingenierie global & 7.4 & Bon niveau, dette d'industrialisation \\
\bottomrule
\end{longtable}

\section{10. Roadmap d'amelioration recommandee}
\subsection{Phase 1 (0-30 jours)}
\begin{enumerate}
  \item Unifier auth/RBAC et conventions roles.
  \item Standardiser schema evenements realtime.
  \item Rendre explicite le mode de donnees affiche (live/sim/fallback).
  \item Supprimer fallback silencieux critiques en production.
\end{enumerate}

\subsection{Phase 2 (30-90 jours)}
\begin{enumerate}
  \item Versionner les contrats API et evenements.
  \item Deplacer simulation vers backend event-driven.
  \item Unifier pipeline IA et instrumenter metriques qualite.
  \item Ajouter audit logs structurants et traçabilite actions.
\end{enumerate}

\subsection{Phase 3 (90-180 jours)}
\begin{enumerate}
  \item Multi-tenant robuste et securite zero-trust inter-services.
  \item Device twin + OTA workflow complet.
  \item Observabilite complete (metrics/logs/traces/SLO).
  \item CI/CD qualite: tests integration, tests contrat, scans securite.
\end{enumerate}

\section{Conclusion}
Smart Shepherd dispose d'une base technique ambitieuse et credible. Le produit est deja fort en UX et demonstration metier. Pour atteindre un standard SaaS IoT/IA enterprise, la priorite n'est pas d'ajouter des ecrans mais de renforcer la coherence transversale: securite, gouvernance des contrats, fiabilite realtime et industrialisation IA.

\vspace{0.5cm}
\textbf{Decision architecturale conseillee:} conserver la stack actuelle, mais lancer un programme d'unification technique sur 3 phases avec jalons de qualite mesurables.

\end{document}
